\chapter{Kanäle}\label{kap:kanaele}

\begin{merksatz}
\boxtitel{Der Leitsatz}
\begin{center}
\renewcommand{\arraystretch}{1.55}
\begin{tabular}{@{}P{2.6cm}P{6.2cm}@{}}
\toprule
\textbf{Nicht} & Speicher durch gemeinsamen Zugriff teilen\\
\textbf{Sondern} & \textbf{Speicher teilen, indem man ihn weiterreicht}\\
\textbf{Der Unterschied} & Ein Wert hat zu jeder Zeit genau einen Besitzer\\
\bottomrule
\end{tabular}
\end{center}

\begin{center}
\begin{minipage}{0.88\textwidth}
\kursiv{Die dritte Zeile ist derselbe Begriff wie in Band 21, Kapitel~3 --- nur wird er hier nicht vereinbart, sondern durch die Bauform nahegelegt. Wer einen Wert über einen Kanal schickt, gibt ihn ab; wer ihn behält und zugleich verschickt, hat einen Wettlauf.}
\end{minipage}
\end{center}

\textbf{Die zwei Arten.}

\begin{center}
\renewcommand{\arraystretch}{1.5}
\begin{tabular}{@{}P{2.4cm}P{2.8cm}P{3.0cm}@{}}
\toprule
& \textbf{Ohne Puffer} & \textbf{Mit Puffer}\\
\midrule
Senden blockiert & \textbf{bis empfangen} & bis Puffer voll\\
Wirkung & \textbf{Gleichschritt} & \textbf{Entkopplung}\\
Gegendruck & \textcolor{gruen}{\textbf{sofort}} & \textcolor{mittel}{verzögert}\\
Verwendung & Übergabe, Signal & \textcolor{gruen}{Lastspitzen}\\
\bottomrule
\end{tabular}
\end{center}

\begin{center}
\kursiv{Die dritte Zeile ist die wichtigste für die Übertragungsschicht: Ein Kanal ohne Puffer bremst den Sender genau dann, wenn der Empfänger nicht nachkommt. Das ist Gegendruck ohne eine einzige Zeile Code --- und die Eigenschaft, an der ein System unter Last steht oder fällt.}
\end{center}

\begin{technik}
\boxtitel{Wie groß der Puffer sein soll}
\begin{center}
\renewcommand{\arraystretch}{1.45}
\begin{tabular}{@{}P{2.6cm}P{2.6cm}P{3.0cm}@{}}
\toprule
\textbf{Größe} & \textbf{Wirkung} & \\
\midrule
$0$ & \textcolor{gruen}{\textbf{voller Gegendruck}} & \textbf{Vorgabe}\\
Klein ($8$--$64$) & \textcolor{gruen}{glättet Schwankungen} & \\
Groß ($10^5$) & \textcolor{akzent}{\textbf{verbirgt das Problem}} & \textcolor{akzent}{Speicher, Verzögerung}\\
Unbegrenzt & \textcolor{akzent}{gibt es nicht} & \textcolor{gruen}{aus gutem Grund}\\
\bottomrule
\end{tabular}
\end{center}

\begin{center}
\begin{minipage}{0.86\textwidth}
\kursiv{Die dritte Zeile beschreibt den häufigsten Fehler: Ein großer Puffer lässt ein überlastetes System eine Weile gesund aussehen und macht den Zusammenbruch später und schwerer. Ein Puffer glättet Schwankungen; er ersetzt keine Kapazität.}
\end{minipage}
\end{center}

\begin{center}
\kursiv{Die letzte Zeile ist eine bewusste Auslassung: Go bietet keinen unbegrenzten Kanal, weil er das Gegendrucksignal zerstört --- man würde den Speicher füllen, statt den Sender zu bremsen.}
\end{center}
\end{technik}
\end{merksatz}

\begin{merksatz}
\boxtitel{\code{select} und die Regeln}
\begin{center}
\begin{minipage}{0.9\textwidth}
\code{select \{}\\
\code{case m := <-eingang:}\\
\code{\ \ \ \ verarbeite(m)}\\
\code{case <-zeit.After(5 * time.Second):}\\
\code{\ \ \ \ return ErrZeitueberschreitung}\\
\code{case <-ctx.Done():}\\
\code{\ \ \ \ return ctx.Err()}\\
\code{\}}
\end{minipage}
\end{center}

\begin{center}
\renewcommand{\arraystretch}{1.5}
\begin{tabular}{@{}P{2.8cm}P{2.4cm}P{2.8cm}@{}}
\toprule
\textbf{Regel} & \textbf{Verhalten} & \\
\midrule
Mehrere bereit & \textbf{zufällig einer} & \textcolor{gruen}{keine Bevorzugung}\\
Keiner bereit & blockiert & \\
Mit \code{default} & \textcolor{akzent}{\textbf{nie blockierend}} & \textcolor{akzent}{Vorsicht}\\
Auf \code{nil}-Kanal & \textbf{nie bereit} & \textcolor{gruen}{\textbf{nützlich}}\\
\bottomrule
\end{tabular}
\end{center}

\begin{center}
\begin{minipage}{0.88\textwidth}
\kursiv{Die erste Zeile verhindert Aushungern: Wenn zwei Fälle bereit sind, wird gewürfelt statt der erste genommen. Die letzte Zeile ist ein Kunstgriff --- einen Kanal auf \code{nil} zu setzen schaltet seinen Zweig ab, ohne den \code{select} umzubauen.}
\end{minipage}
\end{center}

\begin{center}
\kursiv{Der dritte Fall ist die häufigste Falle: Ein \code{select} mit \code{default} in einer Schleife wartet nie und verbraucht einen Kern vollständig --- eine beschäftigte Warteschleife, die keine sein sollte.}
\end{center}

\begin{warnung}
\boxtitelr{Wer schließt den Kanal?}
\begin{center}
\renewcommand{\arraystretch}{1.5}
\begin{tabular}{@{}P{3.0cm}P{2.2cm}P{2.8cm}@{}}
\toprule
\textbf{Handlung} & \textbf{Erlaubt?} & \\
\midrule
Sender schließt & \textcolor{gruen}{\textbf{ja}} & \textbf{die Regel}\\
Empfänger schließt & \textcolor{akzent}{\textbf{nein}} & \textcolor{akzent}{Sender stürzt ab}\\
Zweimal schließen & \textcolor{akzent}{\textbf{nein}} & \code{panic}\\
Auf geschlossenen senden & \textcolor{akzent}{nein} & \code{panic}\\
Von geschlossenem lesen & \textcolor{gruen}{ja} & \textcolor{gruen}{Nullwert, \code{ok=false}}\\
\bottomrule
\end{tabular}
\end{center}

\begin{center}
\begin{minipage}{0.88\textwidth}
\kursiv{Die Asymmetrie ist der Kern: Schließen ist eine Aussage des Senders --- \enquote{es kommt nichts mehr}. Der Empfänger kann diese Aussage nicht treffen. Bei mehreren Sendern braucht es deshalb einen eigenen Mechanismus, meist eine Zählgruppe.}
\end{minipage}
\end{center}
\end{warnung}
\end{merksatz}

\begin{aufgabe}
\begin{enumerate}[label=(\alph*)]
\item Bauen Sie eine Verarbeitungskette.
\item Messen Sie die Wirkung des Gegendrucks.
\item Bestimmen Sie, wann Sperren besser sind.
\end{enumerate}

\loesung
\textbf{(a) Eine Verarbeitungskette.}

\textbf{Die Aufgabe.} Messwerte lesen, umrechnen, prüfen, bündeln, senden.

\begin{center}
\begin{minipage}{0.92\textwidth}
\code{func rechne(ctx context.Context, ein <-chan Rohwert)}\\
\code{\ \ \ \ \ \ \ \ <-chan Messwert \{}\\
\code{\ \ \ \ aus := make(chan Messwert)}\\
\code{\ \ \ \ go func() \{}\\
\code{\ \ \ \ \ \ \ \ defer close(aus)\ \ \ \ \ \ \ \ // Sender schliesst}\\
\code{\ \ \ \ \ \ \ \ for r := range ein \{}\\
\code{\ \ \ \ \ \ \ \ \ \ \ \ select \{}\\
\code{\ \ \ \ \ \ \ \ \ \ \ \ case aus <- umrechne(r):}\\
\code{\ \ \ \ \ \ \ \ \ \ \ \ case <-ctx.Done(): return}\\
\code{\ \ \ \ \ \ \ \ \ \ \ \ \}}\\
\code{\ \ \ \ \ \ \ \ \}}\\
\code{\ \ \ \ \}()}\\
\code{\ \ \ \ return aus}\\
\code{\}}
\end{minipage}
\end{center}

\begin{center}
\begin{minipage}{0.88\textwidth}
\kursiv{Vier Eigenschaften, die zusammengehören: Die Funktion erzeugt ihren Ausgangskanal, schließt ihn selbst, beachtet den Abbruch und gibt einen Nur-Lese-Kanal zurück. Wer alle vier einhält, kann Stufen beliebig verketten, ohne über Lebensdauern nachzudenken.}
\end{minipage}
\end{center}

\textbf{Die Kette.}

\begin{center}
\begin{minipage}{0.86\textwidth}
\code{roh := lies(ctx, geraet)}\\
\code{mess := rechne(ctx, roh)}\\
\code{gut := pruefe(ctx, mess)}\\
\code{pakete := buendle(ctx, gut, 500)}\\
\code{err := sende(ctx, pakete, ziel)}
\end{minipage}
\end{center}

\begin{center}
\renewcommand{\arraystretch}{1.45}
\begin{tabular}{@{}P{3.0cm}P{2.2cm}P{2.8cm}@{}}
\toprule
\textbf{Eigenschaft} & \textbf{Ergebnis} & \\
\midrule
Stufen einzeln prüfbar & \textcolor{gruen}{\textbf{ja}} & \textbf{Kanal hinein, Kanal heraus}\\
Reihenfolge änderbar & \textcolor{gruen}{ja} & \\
Abbruch wirkt überall & \textcolor{gruen}{\textbf{ja}} & \textcolor{gruen}{ein Kontext}\\
Zwischenpuffer & \textcolor{gruen}{je Stufe wählbar} & \\
\bottomrule
\end{tabular}
\end{center}

\begin{center}
\kursiv{Die erste Zeile ist der Prüfgewinn: Jede Stufe ist eine reine Abbildung von Kanal auf Kanal und lässt sich mit einer Liste von Eingaben prüfen --- ohne Gerät, ohne Netz, ohne Attrappe. Das ist die Trennung aus Band 20, Kapitel~6 in nebenläufiger Form.}
\end{center}

\medskip
\textbf{(b) Die Wirkung des Gegendrucks.}

\textbf{Der Versuch.} Der Empfänger wird langsamer als die Sender.

\begin{center}
{\small
\renewcommand{\arraystretch}{1.45}
\begin{tabular}{@{}P{2.4cm}P{2.0cm}P{2.2cm}P{2.2cm}@{}}
\toprule
\textbf{Puffer} & \textbf{Speicher} & \textbf{Verlust} & \textbf{Verzögerung}\\
\midrule
$0$ & \textcolor{gruen}{$\mathbf{2}$ MB} & \textcolor{gruen}{$\mathbf{0}$} & \textcolor{gruen}{$\mathbf{12}$ ms}\\
$1000$ & $6$ MB & $0$ & $310$ ms\\
$10^5$ & \textcolor{akzent}{$140$ MB} & $0$ & \textcolor{akzent}{$\mathbf{31}$ \textbf{s}}\\
$10^6$ & \textcolor{akzent}{$\mathbf{1{,}4}$ \textbf{GB}} & \textcolor{akzent}{\textbf{Absturz}} & \textcolor{akzent}{---}\\
\bottomrule
\end{tabular}}
\end{center}

\begin{center}
\begin{minipage}{0.88\textwidth}
\kursiv{Die dritte Zeile zeigt den eigentlichen Schaden: Kein Verlust, aber einunddreißig Sekunden Verzögerung --- ein Messwert erreicht die Klinik eine halbe Minute zu spät. Bei einer Grenzwertwarnung ist das schlimmer als ein sauberer Verlust mit Meldung.}
\end{minipage}
\end{center}

\begin{center}
\kursiv{Die erste Zeile ist die richtige Wahl: Ohne Puffer bremst der Gegendruck bis zum Gerät durch, und das Gerät kann selbst entscheiden, was es verwirft. Die Entscheidung fällt dort, wo man sie fachlich treffen kann.}
\end{center}

\medskip
\textbf{(c) Wann Sperren besser sind.}

\begin{center}
{\small
\renewcommand{\arraystretch}{1.5}
\begin{tabular}{@{}P{3.4cm}P{2.0cm}P{2.6cm}@{}}
\toprule
\textbf{Lage} & \textbf{Kanal?} & \textbf{Sperre?}\\
\midrule
Werte weiterreichen & \textcolor{gruen}{\textbf{ja}} & nein\\
Ablauf koordinieren & \textcolor{gruen}{ja} & nein\\
\textbf{Zwischen\-speicher schützen} & \textcolor{akzent}{nein} & \textcolor{gruen}{\textbf{ja}}\\
Zähler erhöhen & \textcolor{akzent}{nein} & \textcolor{gruen}{atomar}\\
Konfiguration lesen & nein & \textcolor{gruen}{Lesesperre}\\
\bottomrule
\end{tabular}}
\end{center}

\begin{center}
\begin{minipage}{0.88\textwidth}
\kursiv{Die dritte Zeile widerspricht dem Leitsatz und ist trotzdem richtig: Für einen Zwischenspeicher, auf den viele lesend zugreifen, ist eine Sperre einfacher und schneller als eine Goroutine, die ihn verwaltet. Der Leitsatz ist eine Voreinstellung, kein Verbot.}
\end{minipage}
\end{center}

\textbf{Die Messung.}

\begin{center}
\renewcommand{\arraystretch}{1.45}
\begin{tabular}{@{}P{3.4cm}P{2.0cm}P{2.0cm}P{1.6cm}@{}}
\toprule
\textbf{Vorgang} & \textbf{Kanal} & \textbf{Sperre} & \\
\midrule
Zähler erhöhen & $260$ ns & $\mathbf{18}$ ns & $14\times$\\
Wert übergeben & $\mathbf{95}$ ns & $110$ ns & \\
Zwischenspeicher lesen & $240$ ns & $\mathbf{6}$ ns & $40\times$\\
\bottomrule
\end{tabular}
\end{center}

\begin{center}
\kursiv{Ein Kanal kostet etwa hundert Nanosekunden --- viel für einen Zähler, wenig für eine Netzübertragung. Die Wahl folgt aus der Häufigkeit: Was millionenfach je Sekunde geschieht, braucht keine Nachrichten.}
\end{center}

\begin{warnung}
\boxtitelr{Kanäle lösen keine Entwurfsfehler}
\begin{center}
\renewcommand{\arraystretch}{1.5}
\begin{tabular}{@{}P{3.0cm}P{2.2cm}P{2.8cm}@{}}
\toprule
\textbf{Erwartung} & \textbf{Erfüllt?} & \\
\midrule
Keine Wettläufe mehr & \textcolor{akzent}{\textbf{nein}} & \textbf{wenn man Zeiger schickt}\\
Keine Verklemmungen & \textcolor{akzent}{\textbf{nein}} & \textcolor{akzent}{Kap.~\ref{kap:schiefgeht}}\\
Automatisch schneller & \textcolor{akzent}{nein} & \\
Übersichtlicher Ablauf & \textcolor{gruen}{\textbf{ja}} & \textcolor{gruen}{der eigentliche Gewinn}\\
\bottomrule
\end{tabular}
\end{center}

\begin{center}
\begin{minipage}{0.88\textwidth}
\kursiv{Die erste Zeile ist die verbreitetste Fehlannahme: Wer einen Zeiger über einen Kanal schickt und ihn danach weiter benutzt, hat denselben Wettlauf wie mit gemeinsamem Speicher. Der Kanal überträgt den Besitz nur, wenn man ihn tatsächlich abgibt.}
\end{minipage}
\end{center}
\end{warnung}
\end{aufgabe}
\vspace{-4pt}
\noindent Hoare\zit{13} beschrieb die kommunizierenden Prozesse, aus denen die Kanäle stammen; die Mus\-ter\-vor\-trä\-ge\zit{21} zeigen die wiederkehrenden Bau\-for\-men; das Speichermodell\zit{10} legt fest, welche Sicht\-bar\-keits\-zu\-si\-che\-run\-gen daraus folgen.
